%!TEX TS-program = xelatex
%!TEX encoding = UTF-8
% ------------------------------------------------------------------
% 《__ZH_TITLE__》(中文学习版)
% 改编自 __SOURCE_CITATION__
% 版式: AJbook 文档类 (李文威, CC BY 4.0), 经由 https://github.com/lichuang/latex-template.
% 编译: `make` (即 latexmk, 配置见 latexmkrc: xelatex + biber + xindy).
% ------------------------------------------------------------------
\documentclass[
	draftmark = false,
	fontsetup = font-setup-open.tex,
	titlesetup = titles-setup.tex
]{AJbook}

% 只有一份索引: 不要加 splitindex (xindy 读不了 splitindex 格式, 会得到空索引)
\usepackage[xindy]{imakeidx}
\makeindex[
	columns=2,
	program=xindy,	% 由 latexmk 调用 (见 latexmkrc), 不需要 -shell-escape
	intoc=true,
	options=-M texindy -I xelatex -C utf8,
	title={名词索引}]

\usepackage[unicode, bookmarksnumbered]{hyperref}
\hypersetup{
	pdfauthor = {__AUTHOR__ (原著); 中文学习版},
	pdftitle = {__ZH_TITLE__},
	CJKbookmarks = true}

\addbibresource{references.bib}
% 中文学习版的自定义宏与环境 (在 AJbook 模板之上追加)
\usepackage{listings}
\usepackage{xcolor}
\usepackage{graphicx}
\usepackage{bm}
\usepackage{longtable}

% ---------- 常用记号 (按原书需要增删, 保持与原书记号一致) ----------
\providecommand{\R}{\mathbb{R}}
\providecommand{\Z}{\mathbb{Z}}

% ---------- “补充讲解” 方框 ----------
% 凡是中文版在原文之外新增的讲解/例题/推导, 一律放进 jiedu 环境,
% 与原作者的内容严格分开 (开放许可一般要求标明改动; 对读者也更诚实).
\newtcolorbox{jiedu}[1][]{
	breakable, enhanced,
	colback = teal!4, colframe = teal!60!black, boxrule = 0.3mm,
	coltitle = teal!40!black, colbacktitle = teal!15,
	fonttitle = \sffamily\bfseries,
	title = {补充讲解\ifx\relax#1\relax\else：#1\fi},
	before skip = 0.6em, after skip = 0.6em
}
% 行内的“译注” (改正笔误、说明省略、简短补充)
\newcommand{\yizhu}[1]{\footnote{\textbf{译注：}#1}}
% 课件改编时标注出处页: \slideref{12} -> [slide 12]
\newcommand{\slideref}[1]{\textsf{\small\color{gray}[slide #1]}}

% ---------- 代码 ----------
\lstset{
	language = Python,
	basicstyle = \ttfamily\small,
	keywordstyle = \color{blue!70!black},
	commentstyle = \color{green!40!black},
	stringstyle = \color{red!60!black},
	frame = single, framerule = 0.3pt, rulecolor = \color{gray!60},
	breaklines = true, columns = fullflexible, keepspaces = true,
	showstringspaces = false
}


\renewcommand*\arraystretch{1.3}
\numberwithin{equation}{chapter}

% 需要时只编译部分章节: \includeonly{chapters/ch03}

\begin{document}
	\frontmatter
	% 扉页与前言
\thispagestyle{empty}
\begin{center}
	\vspace*{3cm}
	{\Huge\heiti __ZH_TITLE__}\\[1.2em]
	{\Large 中文学习版}\\[3em]
	{\large 原著：__AUTHOR__}\\[0.5em]
	{\normalsize \emph{__EN_TITLE__}__EDITION__}\\[\fill]
\end{center}
\clearpage

\chapter*{版权与说明}
\addcontentsline{toc}{chapter}{版权与说明}

\paragraph*{原文} __SOURCE_CITATION__. __LICENSE_SENTENCE__

\paragraph*{本版} 本版是原文的中文学习改编版. 所作的改动:
\begin{enumerate}
	\item 叙述与讲解译为简体中文; 专有名词、定理名与章节标题保留原文的英文写法, 以便与原文对照、做作业;
	\item 在原文较为简略之处增加了补充讲解、例题与推导细节. 凡是新增内容, 一律放在标题为 “\textbf{补充讲解}” 的彩色方框中, 或以 “\textbf{译注}” 脚注的形式出现;
	\item 改正了原文的若干笔误, 每处都以译注注明原文写法;
	\item 版式使用 AJbook 文档类 (李文威, CC BY 4.0), 经由 \url{https://github.com/lichuang/latex-template} 获得.
\end{enumerate}
原作者并未审阅或认可本改编版; 其中的任何错误均属改编者.

\chapter*{阅读说明}
\addcontentsline{toc}{chapter}{阅读说明}
\paragraph*{术语} 书中术语一律使用原文的英文写法 (例如 __TERM_EXAMPLES__), 中文只用于叙述和讲解. 书末索引也按英文术语排列.
\paragraph*{补充讲解} 彩色方框中的内容是中文版新增的, 不属于原文; 考试与作业以原文为准.

	\tableofcontents

	\mainmatter
	\chapter{__ORIGINAL_CHAPTER_TITLE__}\label{ch:1}

% 章首: 原文的学习目标 (若有)

\section{__ORIGINAL_SECTION_TITLE__}\label{sec:1.1}

% 正文: 中文叙述, 术语英文. 定义处用 \emph{...} 并加索引:
% 一个 \emph{basic feasible solution}\index{basic feasible solution} 是……

\begin{jiedu}[这一段为什么成立]
	新增讲解放在这里: 补全跳过的步骤, 给小的数值例子 (用代码验证过的数字), 解释直觉.
\end{jiedu}

\section{Exercises}
\begin{Exercises}
	\item\label{exo:1.1} \textbf{(原题标题)} 题目译文.\yizhu{提示: 只给方向, 不给完整解答.}
\end{Exercises}

	% \include{chapters/ch02} ...

	\backmatter
	\printbibliography[heading=bibintoc, title={参考文献}]
	{\footnotesize
	\indexprologue{术语与原文一致, 使用英文, 按字母排序.}
	\printindex}
\end{document}
